\documentclass[hyphens,aspectratio=169,dvipsnames]{beamer}
\usepackage{graphicx}
\usepackage{xcolor}
\usepackage{amssymb}
\usepackage{amsmath}
\usepackage{mathtools}
\usepackage{textcomp}
\usepackage{moresize}
\usepackage{framed}
\usepackage{minted}
\usepackage{relsize}
\usepackage{tikz}
\usepackage{comment}
\usetikzlibrary{shapes.geometric, arrows, positioning}
\usetheme{Berlin}
\usecolortheme[RGB={0,95,47}]{structure}
\beamertemplatenavigationsymbolsempty

\makeatletter
\AtBeginEnvironment{minted}{\dontdofcolorbox}
\def\dontdofcolorbox{\renewcommand\fcolorbox[4][]{##4}}
\makeatother

\newcommand{\textpf}[1]{\texttt{\color{black}\fcolorbox{lightgray}{lightgray}{#1}}}

\begin{document}

\begin{frame}[fragile]{Logistics}
    \begin{itemize}
        \item The exam is on Tuesday night, 6:00-8:00. If you can't make this time, let us know ASAP.
        \item Topics covered will be
            \begin{itemize}
                \item integers
                \item structs
                \item arrays
                \item pointers
                \item \textpf{argc} and \textpf{argv}
                \item syscalls
                \item binary and hex
                \item 2's complement
                \item sign-extension
            \end{itemize}
        \item Quizzes are now graded for completion.
    \end{itemize}
\end{frame}

\begin{frame}[fragile]{\textpf{const}}
    Example!
\end{frame}

\begin{frame}[fragile]{\textpf{printf}}
    Example!
\end{frame}

\begin{frame}[fragile]{Activity}
    Implement a function \textpf{atoul}, which takes 2 arguments, a string \textpf{s} and an \textpf{unsigned char b} between 1 and 16, and returns the \textpf{unsigned long} obtained by parsing \textpf{s} in base \textpf{b}.

    For example, \textpf{atoul("123", 10)} should return \textpf{123}, and \textpf{atoul("100", 2)} should return 4.
\end{frame}

\end{document}
